% GUIDA VELOCE:
% --------------------------------------------------------------------
% X INIZIA UN UNOVO CAPITOLO:
% \chapter{??? NOME CAPITOLO}}
% \section{ ?? }
% \subsection{ ?? }
%
% --------------------------------------------------------------------
% PAROLA CONTENUTA NEL GLOSSARIO:
% scrivere la parola seguita da $^g$
% esempio: User$^g$
%
% --------------------------------------------------------------------
% PER ANDARE A CAPO SENZA RIENTRO INSERIRE:
% \\
%
% --------------------------------------------------------------------

% GRASSETO:
% \textbf{parola}
%
% --------------------------------------------------------------------
% CORSIVO:
% \emph{parola}
% --------------------------------------------------------------------
% PER SCRIVERE IN ROSSO:
% \red{parola}
%
% --------------------------------------------------------------------
% PER SCRIVERE TRA VIRGOLETTE
% ''parola''
%
% --------------------------------------------------------------------
% PER EVITARE IL RIENTRO AUTOMATICO DI UN CAPOVERSO:
% \noindent testo....
%
% --------------------------------------------------------------------

% PER SCRIVERE CARATTERI PARTICOLARI COME: { } _ ecc.. SCRIVERLI PRECEDUTI DA \
% ES: \{ \_
%
% --------------------------------------------------------------------
% X INSERIRE UN LINK:
% \url{http://www.math.unipd.it/~tullio/IS-1/2011/Progetto/C3.pdf}
%
% --------------------------------------------------------------------
% PER COMMENTARE INTERE PARTI:
% \comment{ comment }
%
% --------------------------------------------------------------------
% PER SCRIVERE NOTE DURANTE IL TESTO:
% parola \footnote{ note riguardanti la parola }
%
% --------------------------------------------------------------------
% PER SCRIVERE CODICE SORGENTE:
%
% \lstset{language=c++,
% stringstyle=\color{blue}\textrm,
% commentstyle=\rmfamily, numbers= none}

% \begin{lstlisting}
% CODICE
% \end{lstlisting}
%
% --------------------------------------------------------------------
% !!!!!!!! PER COSE + COMPLESSE VEDI: !!!!!!!!!!!!!!!!!!!!!!!
% !!!!!!!! PMAC/latex/GUIDA LATEX!!!.tex !!!!!!!!!!!!!!!!!!!!!!!

% per tutto il resto chiedi a lory prima di fare/scrivere cazzate !!!!!!!!!!



\documentclass[10pt,a4paper]{article}

\usepackage[italian]{babel}
\usepackage[T1]{fontenc}
\usepackage[utf8x]{inputenc} % uso utf8x xk x linux, mentre latin1 è per windows
\usepackage{lmodern} %insieme di font molto completo consigliato da LatexFacile pg13 in basso
\usepackage{microtype} %migliora riempimento delle righe. vedi LatexImpaziente pg41
%attiva il rientro di ogni prima riga di ogni sezione: capitolo,paragrafo ecc. vd LatexImpaziente pg41
\usepackage{indentfirst}
\usepackage{graphicx} % per inseire immagini
\usepackage[usenames,dvipsnames]{color}
\usepackage{lastpage} %serve per poter scrivere page 1 of N
% setta i bordi della pagina: dx e sx 3.2cm di rientro + nel lato di rilagatura rientra di altri 0mm
\usepackage[a4paper,top=2.7cm,bottom=2.7cm,left=2.7cm,right=2.7cm, bindingoffset=0mm]{geometry}
\usepackage{listings} % per inserire codice sorgente
\usepackage{float} % per gestire oggetti flottanti ( es immagini tabelle posizionebili con "H" che forza il posizionamento nel punto specifico )

% serve per creare tabelle lunghe + di una pagina con \begin{longtable} (vd Tabelle.pdf pg11-12)
\usepackage{longtable}

\usepackage{fancyhdr} % per impostare lo stile della pagina più personalizzato, + fancyhdr ( per regolare testatina e piè di pagina ) vedi itfancyhrd


\pagestyle{fancy}
% settaggi di pagestyle(fancy)
\lhead{\includegraphics[scale=0.20]{images/nastro}}
%\chead{}
\rhead{\textbf{Tecnologie Web \\ Data: \DataRilascio}}
\lfoot{\NomeDocumento}
\cfoot{}
\rfoot{ \textbf \thepage\ of \pageref{LastPage}}
\renewcommand{\footrulewidth}{0.4pt}

%ridefinisco il plain per cosare l'indice (a questo punto si potrebbe lasciare tutto il documento in plain
\fancypagestyle{plain}{
%\lhead{\includegraphics[scale=0.20]{SevenFold_small}}
%\chead{}
\rhead{\textbf{{ Tecnologie Web}}}
\lfoot{\NomeDocumento}
\cfoot{}
\rfoot{ \textbf \thepage\ of \pageref{LastPage}}
\renewcommand{\footrulewidth}{0.4pt}
}

% da ultimo:
\usepackage{hyperref} %x l'interpretazione di indirizzi o link ipertestuali (vd LatexImpaziente pg47 )
\hypersetup{backref, colorlinks=true, linkcolor=black, urlcolor=black}

\usepackage{url} % x l'interpretazioni di internet o link ipertestuali (vd LatexImpaziente pg47 )
%\UrlFont{color =blue}
%\urlstyle{helvetic}

% Define a new 'leo' style for the package that will use a smaller font.
\makeatletter
\def\url@leostyle{%
  \@ifundefined{selectfont}{\def\UrlFont{\sf}}{\def\UrlFont{\small\ttfamily}}}
\makeatother
%% Now actually use the newly defined style.
\urlstyle{leo}


\newcommand{\mail}[1]{\textcolor{Black}{ \texttt{#1}}} %per interpretare mail (vd LatexImpaziente pg47 )
\newcommand{\cambiaFont}[2]{{\fontencoding{T1}\fontfamily{#1}\selectfont#2}}
\newcommand{\red}[1]{ \textcolor{red}{#1} } % per scrivere testo in rosso
\newcommand{\blue}[1]{ \textcolor{blue}{#1} } % per scrivere testo in blu
\newcommand{\comment}[1]{} % per inserire commenti

\newcommand{\attribute}[2]{ \item[\textcolor{PineGreen}{ \texttt{#1}}] \textcolor{PineGreen}{\texttt{#2\\}}\ \ \ }
\newcommand{\method}[2]{ \item[\textcolor{MidnightBlue}{ \texttt{#1}}] \textcolor{MidnightBlue}{ \texttt{#2\\}}\ \ \ }

\newcommand{ \class}[1]{ \item[ \textcolor{MidnightBlue}{\texttt{-}}] \textcolor{MidnightBlue}{\texttt{#1}}}
\newcommand{ \classe}[1]{ \noindent \textcolor{MidnightBlue}{\texttt{-}} \textcolor{MidnightBlue}{\texttt{#1}}}

\newcommand{\tab}{\hspace*{1.5em}}


% INSERIRE QUI IL NOME DEL DOCUMENTO SEGUITO DA UNO SPAZIO
% ( così il nome si imposta in automatico nelle varie ricorrenze standard)
\newcommand{\NomeDocumento}{Relazione progetto: \textit{"AllStreaming"}}

% INSERIRE QUI LA DATA DEL RILASCIO DELLA VERSIONE ATTUALE
\newcommand{\DataRilascio}{19-06-2012}

\newcommand{\serie}{\textit{Serie Tv}}
\newcommand{\film}{\textit{Film}}


\begin{document}


% --------------------------------------------------------------------

% TITOLO ( 1° pagina)

\vspace*{1cm}
\begin{center}

%\cambiaFont{Cyklop}{Sevenfold}
%\cambiaFont{fve}{\Huge{Sevenfold}}


\vspace{1.5cm}

\textit{\Huge{Relazione per il progetto \\del corso di Tecnologie Web:}}\\
\vspace{1.3cm}
\cambiaFont{fve}{\Huge{\color{blue}{""AllStreaming""}}}
\vspace*{1.8cm}

\large{Università degli studi di Padova\\
anno accademico 2011-2012}

\vspace*{1cm}

\includegraphics[scale=0.35]{images/nastro10e.png}

\end{center}


% --------------------------------------------------------------------

% INFORMAZIONI DEL DOCUMENTO ( 1° pagina)

\vspace*{2cm}
\begin{center}

\begin{tabular}{ r | l }
\multicolumn{2}{c}{\textbf{\Large{Informazioni del documento}} }\\
\hline
\rule[-1.5mm]{0mm}{0.7cm}
\textbf{Titolo del documento} & \NomeDocumento\\
\rule[-1.5mm]{0mm}{0.5cm}
\textbf{Date}& \DataRilascio\\
\rule[-1.5mm]{0mm}{0.5cm}
\textbf{Autori}& \underline{Lorigiola Giacomo} [ 592992 ]\\
& \underline{Spiga Gabriele} [592141]\\
& \underline{Luca Guerra} [???]\\
\rule[-1.5mm]{0mm}{0.5cm}
\textbf{Destinatario}&Prof.essa Ombretta Gaggi\\
\end{tabular}

\end{center}

% --------------------------------------------------------------------

% SOMMARIO ( 2° pagina)

\newpage

\vspace*{0.5cm} % il vertical space va preceduto da una riga vuota!!!
\begin{center}

\textbf{{\Large{Sommario}}}

Questo documento intende presentare una analisi al progetto di Tecnologie Web intitolato "AllStreaming".


\vspace*{0.2cm} % il vertical space va preceduto da una riga vuota!!!

\end{center}



% --------------------------------------------------------------------
% INDICI:

\newpage

% INDICE CAPITOLI
\tableofcontents % genera l'indice di tutto il documento

\let\cleardoublepage\clearpage % toglie la pagina bianca dopo l'indice

% INDICE TABELLE
% \listoftables

% INDICE FIGURE
\listoffigures


% --------------------------------------------------------------------

% INTRODUZIONE ( 1° CAPITOLO ) QUESTO CAPITOLO VA MESSO IN OGNI DOCUMENTO!!!!!!!!

\newpage


\section{Introduzione}
L'obbiettivo del sito è fornire un portale unitario per la visualizzazione di link attraverso i quali poter accedere alla visualizzazione in streaming di contenuti video come film o serie Tv.
Gli utenti potranno selezionare il video di loro interesse scegliendo tra quelli disponibili nelle apposite sezioni \film \ e \serie, e seguendo i link riportati, accedere ai più famosi portali di streaming per la visualizzazione del video scelto.\\


Il sito realizzato è stato caricato nel server di tecnologie-web della facoltà, alle login \textbf{glorigio} e \textbf{lguerra}.




\section{Mappa del sito}
suddivisione pagine\\
schema logico/mappa\\

home\\
serie Tv -> serie\\
film -> film\\
Registrazione / Profilo \\

aboutUS\\

parlare di ogni pagina cosa fa/serve\\


\section{Struttura delle pagine (layout)}
layout normale (con immagine)\\
layout smartphone (con immagine)\\
layout stampa (con immagine)\\


suddivisione pagina:\\
header\\
right(nav) center left(ultimi commentati, novità, utenti iscritti)\\
footer\\



\section{Tecnologie utilizzate e "passed"}

xhtl css js perl




\subsubsection{Client class }
\begin{center}
\begin{figure}[H]
\centering
%\includegraphics[scale=0.45]{clientClassDiagramm.png}
\caption{Client class diagram}
\label{ww-lists }
\end{figure}
\end{center}



\subsubsection{Server class }
\begin{center}
\begin{figure}[H]
\centering
%\includegraphics[scale=0.5]{serverClassDiagramm.png}
\caption{Server class diagram}
\label{ww-lists }
\end{figure}
\end{center}




\end{document}
